\scanpage{24}

\uline{Sous-espace \deleted{\unsure{algébr.|algébrique}} fermé de $\mathbb E_k^I$} \uline{espace alg. affine}

\[
\text{\textcircled{4}}\quad V_S=V_J
\]

\inserted{idéal engendré, c'est le syst. d'équations le plus naturel}

\margin{référence au ``th. de la base'' de Hilbert ($k$ corps \deleted{\unsure{plus}} gén, noeth) $\Rightarrow k[T_1-\ T_n]$ noeth}

\[
A_J=P_I/J
\]

\diagram{\begin{tikzcd}[ampersand replacement=\&,row sep=small,column sep=small]
k^{\prime I}\arrow[r,phantom,"\simeq"]\arrow[d,phantom,"\rotatebox{90}{$\subset$}"]\&\mathrm{Hom}_{k\text{-}\mathrm{alg}}(P_I,k')\arrow[d,phantom,"\rotatebox{90}{$\subset$}"]\\
V_J(k')\arrow[r,phantom,"\simeq"]\&\{\deleted{\unsure{\mathrm{hom}}}u:P_I\to k'\mid\deleted{\unsure{u\ \text{nul sur }J}}\inserted{u(J)=0}\ (\inserted{\text{$u$ nul}}\ \text{sur les }F_j)\}\arrow[d,phantom,"\rotatebox{90}{$\simeq$}"]\\
\&\mathrm{Hom}_k(A_J,k')
\end{tikzcd}}

\deleted{\unsure{4}} Les foncteurs de la forme $V_S$ ($I,S$ convenables) sont, à isom près, les foncteurs représentables $V_A$ donner un isom. entre $V_A$ et un $V_S$ \deleted{\unsure{rev}} revient au $=$ que de se donner un isom $A\simeq P_I/J$ ($J$ idéal de $P_I$ engendré par $S$), \deleted{\unsure{donc}} ; se donner un plongement $V_A\subset\mathbb E_k^I$ \deleted{\unsure{revient à se donner un}} \inserted{identifiant $V_A$ à un sous-esp. alg fermé de $\mathbb E_k^I$}

% slanted
\margin{utilisant le \unsure{résultat sur les|résult. sur les|notion de} foncteurs représentables}

[Plus généralement, tout plongement $V_A\hookrightarrow\deleted{\unsure{E}}\mathbb E^I$ correspond à un épimorphisme d'anneaux $P_I\to A$ (mais un tel épimorphisme n'est pas en général surjectif \ldots] s'identifient aux hom \deleted{\unsure{surjectifs}} $P_I\to A$, ou encore aux systèmes de gén. $(a_i)_{i\in I}$ de $A$ comme \uline{algèbre} sur $k$.

\margin{Rem}

\textcircled{5} \uline{Cor} La cat. des esp. alg. affines sur $k$ est anti-équivalente à celle des algèbres sur $A$. Ceux qui se plongent comme sous-espaces alg. fermés dans quelque $\mathbb E_k^n$ ($n$ entier $\geq0$ fini) \deleted{sont ceux} correspondent aux $A$ qui \deleted{\unsure{sont un syst. de générateurs finis|sont des $k$-alg. de type fini}} \inserted{sont des $k$-alg. de type fini}.

Les \deleted{\unsure{idéaux de}} sous-espaces algébriques de $\mathbb E_k^I$ correspondent aux idéaux de $P_I$. \deleted{\unsure{en}} $\bigl[$ \uline{\unsure{si on connait|si on se donne|$J$ connu quand on} $=$ connait} $V_J$ \ldots \deleted{\unsure{$J'$}}$J$

\[
\begin{aligned}
V_J\subset V_{J'}&\Longleftrightarrow J'\subset J\\[-2pt]
&\Leftarrow\left\{\begin{aligned}&=\{f\in P_I\mid f(x)=0\ \ \forall\\&\qquad k'/k\text{ et }x\in V_J(k')\}\end{aligned}\right.
\end{aligned}
\]

% archive sheet number: 23
